\documentclass[12pt]{article}
\usepackage{amsmath}
\usepackage{amssymb}
\usepackage{amsthm}
\usepackage{thmtools}
\declaretheorem[numberwithin=section]
{theorem, definition}
\declaretheorem{lemma, proposition, corollary, examples}[
style=plain,
numberwithin=theorem
]

\author{Shan Gao}
\title{Geo topo 1 notes}
\date{fall 2023}
\begin{document}
\maketitle
\begin{abstract}
	Personal notes of MATH 576 taught by Professor Joel Kamnitzer during the fall semester of 2023 
\end{abstract}
\section{Lecture 1}
\begin{definition}[Topological space]\label{def:ts}
    A topological space is a set $X$ along with a collection of subsets which are called open set
    \begin{enumerate}
        \item $\emptyset$, $X$ are open 
        \item let 
            $$
            \{U_\alpha\}_{\alpha \in I}
            $$ be any collection of open sets, where $I$ is an arbitrary index set (not necessarily countable), then
            $$
            \bigcup_{\alpha \in I} U_\alpha
            $$ is open
        \item if $U_1$, $U_2$ are open sets then $U_1 \cap U_2$ is open (open under finite intersections)
    \end{enumerate}
\end{definition}
\begin{examples}
    $\mathbb{R}^n$ with it's usual open sets : $U \subset \mathbb{R}^n$ is open if for any $x \in U$, there exists a $r > 0$ such that $B_r(x) \subseteq U$
\end{examples}
\begin{lemma}[topology on $\mathbb{R}^n$]\label{lem:rntopo}
    We can see that this defines a topology on $\mathbb{R}^n$, we just need to check some of the conditions. The first and second conditions are trivial to verify. As for the third condition, letting $U_1$ and $U_2$ be any open set, we know that, for each $x \in U_1 \cap U_2$, there are suitable radii $r_1, r_2$ that we can use to fit an open ball in each one of the latter sets, we simply need to take $r = \min(r_1,r_2)$, and we are done.
\end{lemma}
\begin{definition}[metric space]
    A metric space $(X,d)$ is a set endowed with a distance map $d \colon X \times X \to \mathbb{R}$, such that         \begin{itemize}
        \item $d(x,y) > 0$ and $d(x,y) = 0 \implies x = y$
        \item $d(x,y) = d(y,x) \forall x,y \in X$
        \item $d(x,z) \leq d(x,y) + d(y,z) \forall x,y,z \in X$
    \end{itemize} 
\end{definition}
\begin{lemma}[any metric space is a topology]
    Any metric space is a topology using the following rule for open sets: Any set $U \in X$ is open if $\exists r > 0$ such that $B_r(x) = \{y \in X \colon d(y,x) < r\} \subseteq U$, by lemma \ref{lem:rntopo}
\end{lemma}
This does not mean that there is a bijection between metric spaces and topological spaces. \textbf{many metrics produce the same topology},for example, $\mathbb{R}^n$ with the $l_p$ metric generate the same topology for any $p$
\begin{examples}
    \begin{itemize}
        \item $X = \{a,b\}$ the open sets are $\emptyset, X$ and any of the two elements, $\{a\}$. Easy to verify that this is a topology.
        \item "Trivial topology", any set $X$ can be made a toplogy with the open sets being only $\emptyset$ and $X$ itself. This is the same topology generate by the metric $d(x,y) = 0 \forall x,y \in X$. 
        \item  Every subset of $X$ is open, so open sets $ = 2^x$, this is the same as setting the distance between two elements to be $0$ if they are the same and $1$ if they are distinct, that way we can do something like $B_{\frac{1}{2}} = \{x\}$, so every subset can be open. 
        \item let $X$ be any set, we can define the \textbf{cofinite topology} : $U \in X$ is open, if $U^c$ is finite, and $\emptyset$ is open. Easy to check that this is a topology.
    \end{itemize}
\end{examples}
\begin{definition}[closed set definition for topology]
    We can define a topology using closed sets, $F \subseteq X$ is called closed if $F^c = X\setminus F$ is open. These closed sets satisfy these conditions
    \begin{itemize}
        \item $\emptyset, X$ are closed
        \item $\{F_\alpha\}_{\alpha \in I}$ are closed then $ \bigcup_{\alpha \in I} F_\alpha$ is closed.
        \item if $F_1, F_2$ are closed, then $F_1 \cup F_2$ is closed. 
    \end{itemize}

\end{definition} 
\begin{examples}
    We supplement this with an example of a cofinite topology, letting $K$ be a field and $f \in K[z]$, we denote the vanishing set of $f$ as $V(f) = \{a \in K \colon f(a) = 0\}$, for example if $K = \mathbb{R}$ and $f = (x-1)(x-2)$, then $V(f) = \{1,2\}$. $\{ V(f) \colon f \in K[z]\}$ is the collection of all finite sets, if you set these sets as the closed sets, then you recover the cofinite topology.
\end{examples}
\begin{definition}[sequential limits]
    $X$ is a topological space and $\{x_n\}_{n \geq 1}$ is a sequence in $X$ (basically a map from $\mathbb{N} \to X$. We say that $x \in X$ is a limit of the sequence if, for all open neighborhoods $U$ (a open subset $U$ such that $x \in U$) of $x$, $\exists N$ such that if $n \geq N$, then $x_n \in U$.  
\end{definition}
\begin{definition}[continuous]
    $X,Y$ are two topological spaces, then $f$ is continuous if $\forall U \subseteq Y$ open, $f^{-1}(U)$ is open in $X$ 
\end{definition} 
\begin{lemma}[sequential continuity]\label{lem:seqcont}
	let $f \colon X \to Y$ be a continuous function, and $(x_n) \subseteq X$ is a sequence in $X$, if $x_n \to x \in X$, then $f(x_n) \to f(x)$. 	
\end{lemma}
\begin{proof}
	We know that $x_n \to x$ in $X$, we want to show that $f(x_n) \to f(x)$ if $f$ is continuous. Precisely, we want to show that any open neighborhood of $f(x)$ contains \textbf{all but finitely many} $f(x_n)$s. Let $U \subseteq Y$ be open open neighborhood of $f(x)$, since $f(x) \in U$, then $x \in f^{-1}(U)$. The latter set is open, by continuity of $f$, thus there exists an $N$ such that $\forall n \geq N$, $x_n \in f^{-1}(U)$, therefore, $f(x_n) \in f(f^{-1}(U)) \subseteq U$. 
	
\end{proof}
There are some warnings that we need to consider: 
\begin{itemize}
	\item A sequence might have more than 1 limit. For example,  with the trivial topology, any sequence converges to every other point, since every point has only one open neighborhood, which is the whole set. 
	\item The converse of lemma \ref{lem:seqcont} might not be true. 
\end{itemize}
\begin{definition}
	$f\colon X \to Y$ is called a homeomorphism if 
	\begin{itemize}
		\item $f$ is continuous
		\item $f$ is a bijection
		\item $f^{-1} \colon Y \to X$ is continuous
	\end{itemize}
	In this case, we can write that $X \cong Y$, such that $f(f^{-1}(x)) = f^{-1}(f(x)) = id$.
	\textbf{warning} : First and second condition does not necessarily imply the third.
\end{definition}
\begin{examples}[Stuff with Zhao's cool example]
asd		
\end{examples}
\end{document}
